\section{وقفه به‌عنوان ابزار حمله}

\textbf{بخش اول: حملهٔ \lr{Heckler}}

\subsection*{بخش آ – چرا گردانندهٔ فراخوانی سیستمی در اجرای ناخواسته به‌صورت وقفه خطرناک است؟}

تله و وقفه از یک جدول بردار و یک مسیر مشترک ورود به هسته استفاده می‌کنند. گردانندهٔ فراخوانی سیستمی بر چند فرض بنا شده است:
\begin{itemize}
  \item کسی با اجرای عمدی دستور تله (مثل \code{int 0x80}) آن را صدا زده است و نشانی بازگشت، حالت و ثبات‌ها بر اساس قرارداد فراخوانی آماده شده‌اند.
  \item هسته «دستور وقفه‌دهنده» را بررسی می‌کند تا نوع را بفهمد و سپس پارامترهای ثبات را کنترل و فراخوانی را اجرا می‌کند؛ یعنی منبع و نیت فراخوانی را مفروض می‌گیرد.
  \item فراخواننده یک پردازهٔ کاربر است که خودش از اثرهای فراخوانی (فقط نتیجه در یک ثبات) آگاه است.
\end{itemize}

وقتی همین گردانندهٔ همین بردار با یک وقفهٔ تزریقی، یعنی ناهم‌زمان و در یک لحظهٔ دلخواه، اجرا شود، هیچ‌کدام از این فرض‌ها برقرار نیست. سخت‌افزار برای وقفه‌های خارجی (و تزریق‌شده) بررسی سطح دسترسی دروازه را انجام نمی‌دهد (این بررسی فقط برای دستور نرم‌افزاری \code{int n} است)، و مسیر مشترک هم نمی‌تواند تشخیص دهد که وقفه از یک دستور واقعی آمده است یا از بیرون. در نتیجه هسته محتوای تصادفی ثبات‌های پردازهٔ قربانی را به‌عنوان شمارهٔ فراخوانی و آرگومان می‌پذیرد، یک فراخوانی سیستمی دلخواه را با اختیارات پردازه اجرا می‌کند و نتیجه را روی ثبات‌هایی می‌نویسد که پردازه در حال استفاده از آن‌هاست. مهاجم زمان تزریق را طوری انتخاب می‌کند که اجرای پردازه را به نفع خود تغییر دهد.

\subsection*{بخش ب – اهمیت مقدار ثبات‌ها در لحظهٔ حمله}

هسته شمارهٔ فراخوانی را از ثبات \code{rax} و آرگومان‌ها را از ثبات‌های \code{rdi}، \code{rsi}، \code{rdx}، \code{r10}، \code{r8} و \code{r9} می‌خواند و مقدار برگشتی را هم در \code{rax} می‌گذارد. نتیجه:
\begin{itemize}
  \item مقدار ثبات‌ها در لحظهٔ تزریق از نظر مهاجم «ورودیِ فراخوانی» است. مهاجم با انتخاب لحظه می‌تواند شمارهٔ فراخوانی و آرگومان‌های مطلوب را بسازد، چون ثبات‌ها محتوای محاسبه‌ای را دارند که پردازه در همان لحظه انجام می‌دهد.
  \item نتیجهٔ فراخوانی ساختگی روی \code{rax} می‌نشیند. اگر \code{rax} در آن لحظه مقدار برگشتی تابعی مثل بررسی گذرواژه باشد (تابع در \code{rax} نتیجه می‌گذارد)، مقدار برگشتی با نتیجهٔ فراخوانی ساختگی (مثلاً صفر یعنی موفقیت) جایگزین می‌شود و شاخهٔ برنامه عوض می‌شود. روش دور زدن احراز هویت در \lr{sudo} و \lr{OpenSSH} همین است.
  \item چون سیستم‌عامل مهمان محرمانه است، ناظر نمی‌تواند حافظه و ثبات‌ها را ببیند. او حالت ثبات‌ها را با زمان‌بندی (مثلاً شمارش وقفه‌ها یا رویدادها پیش از ورود به تابع حساس) حدس می‌زند و بنابراین حمله به زمان‌بندی دقیق و مقدار ثبات‌ها وابسته است.
\end{itemize}

\subsection*{بخش ج – ویژگی روال رسیدگی امن و یک راه دفاع}

روالی در هر لحظه بی‌خطر است که (۱) وضعیت‌ ناشی از اجرای خودش را به‌طور کامل ذخیره و بازیابی کند و هیچ اثری روی برنامهٔ وقفه‌خورده نگذارد؛ (۲) فرضی دربارهٔ محتوای ثبات‌ها یا زمان فراخوانی نداشته باشد و ورودی‌هایش را اعتبارسنجی کند؛ (۳) تنها از مسیر و منبع مجاز قابل ورود باشد؛ یعنی بداند به‌دلیل یک دستور واقعی در حالت کاربر فراخوانده شده است؛ (۴) بازورود‌پذیر باشد و با ماسک‌کردن مناسب وقفه‌ها به وضعیت ناسازگار نرسد.

\textbf{راه دفاع پیشنهادی:} چند لایه ممکن است:
\begin{itemize}
  \item \textbf{سطح هسته:} اگر مسیر \code{int 0x80} (فراخوانی سیستمی قدیمی ۳۲‌بیتی) لازم نیست، غیرفعال شود (برای مثال با تنظیمات \code{ia32\_emulation}) تا بردار \code{0x80} گردانندهٔ فعال نداشته باشد. لینوکس برای مهمان‌های \lr{Intel TDX} و \lr{AMD SEV} این قابلیت را به‌صورت پیش‌فرض غیرفعال می‌کند.
  \item \textbf{سطح سخت‌افزار و ناظر:} محدودکردن بردارهایی که ناظر اجازه دارد تزریق کند (مثل ویژگی تزریق محدود/جایگزین در \lr{AMD SEV-SNP} و بررسی تزریق در \lr{Intel TDX})، تا بردارهای ویژه‌ای مثل \code{0x80} و بردارهای استثنا قابل‌تزریق نباشند.
  \item \textbf{سطح گرداننده:} اجرای گردانندهٔ فراخوانی سیستمی فقط وقتی مجاز باشد که با دستور واقعی در حالت کاربر آمده باشد (برای نمونه ثبت منشأ ورود در مسیر مشترک و رد وقفه‌های بدون منشأ معتبر).
\end{itemize}

\textbf{بخش دوم: نشت اطلاعات از طریق وقفه‌ها (حملهٔ شمارنده در مرورگر)}

\subsection*{بخش د – زبانه‌ها چگونه از هم باخبر می‌شوند و چرا هر وب‌گاه الگوی خودش را دارد؟}

دو زبانه حافظهٔ مشترک ندارند، ولی \textbf{سخت‌افزار و هسته را به اشتراک گذاشته‌اند}. زبانهٔ مهاجم حلقه‌ای می‌چرخاند که یک شمارنده را زیاد می‌کند و در هر بازهٔ زمانی کوتاه (مثلاً چند میلی‌ثانیه) مقدار آن را ثبت می‌کند. هر وقت \lr{CPU} (یا همان هستهٔ پردازنده‌ای که حلقه روی آن است) به کار دیگری مشغول شود، شمارنده کمتر زیاد می‌شود و در دنباله نقطهٔ افتادگی می‌ماند. دو عامل اصلی این افتادگی‌ها هستند:
\begin{itemize}
  \item \textbf{وقفه‌ها:} هر وقفه (کارت شبکه، پردازندهٔ گرافیکی، دیسک، زمان‌سنج، وقفهٔ بین‌هسته‌ای) \lr{CPU} را متوقف می‌کند تا گرداننده اجرا شود؛ پردازنده با هر وقفه وضعیت را ذخیره می‌کند و به گردانندهٔ مربوط می‌رود. مدت و زمان این وقفه‌ها از رفتار «زبانهٔ قربانی» می‌آید.
  \item \textbf{رقابت بر منابع مشترک:} کش، پهنای باند حافظه، منابع داخلی هسته و فرکانس پردازنده.
\end{itemize}

هر وب‌گاه الگوی خاص خودش را می‌سازد، چون اندازه، تعداد و ترتیب منابع و درخواست‌های شبکه، زمان‌بندی دریافت بسته‌ها، رندر و اجرای جاوااسکریپت (و در نتیجه فعالیت \lr{GPU} و \lr{CPU}) برای هر وب‌گاه متفاوت است. نتیجه یک «اثر انگشت» زمانی از وقفه‌ها و بار \lr{CPU} است که شمارنده آن را به‌صورت یک سری زمانی ثبت می‌کند و مدل یادگیری ماشین (مثل شبکهٔ عصبی کانولوشنی) آن را به نام وب‌گاه نگاشت می‌کند.

\subsection*{بخش ه – سنجش سه روش دفاعی}

\textbf{۱) فرستادن وقفه‌های دستگاه به یک هستهٔ دیگر.} وقفه‌های دستگاه‌ها (شبکه، \lr{GPU}، دیسک) دیگر هستهٔ مهاجم را متوقف نمی‌کنند؛ پس بزرگ‌ترین بخش سیگنال از بین می‌رود. اما کامل نیست: وقفهٔ زمان‌سنج و وقفه‌های بین‌هسته‌ای (در ادامه) باقی می‌مانند؛ کار نرم‌افزاری مرتبط با وقفه‌ها (مثلاً پردازش بسته‌ها) ممکن است هنوز روی هستهٔ مهاجم اجرا شود؛ و کانال‌های دیگری مثل کش و پهنای باند حافظه باقی‌اند. هزینه: یک هستهٔ پرمشغله برای همهٔ وقفه‌ها (افت کارایی و تأخیر) و نیاز به چسباندن پردازه‌ها به هسته. \textbf{میزان جلوگیری: متوسط (کاهش قابل توجه ولی نه حذف کامل).}

\textbf{۲) کم‌دقت و نوفه‌دار کردن ساعت.} حمله به ساعت دقیق نیاز ندارد: شمارنده فقط تعداد تکرار حلقه در یک بازهٔ چندمیلی‌ثانیه‌ای را می‌شمارد و وقفه‌ها (از چند میکروثانیه تا میلی‌ثانیه) در این مقیاس ظاهر می‌شوند. یعنی ساعتی با دقت چند ده یا صد میکروثانیه هم هنوز همین کار را می‌کند. مهاجم می‌تواند ساعت جایگزین هم بسازد (با نخ شمارنده و \lr{SharedArrayBuffer}). این روش فقط وقتی اثر دارد که ساعت چنان خراب شود که مرورگر عملاً بی‌فایده شود. \textbf{میزان جلوگیری: کم (ضعیف)}. این روش بیشتر در برابر حمله‌هایی کارساز است که به ساعت دقیق نیاز دارند.

\textbf{۳) تزریق وقفه‌های ساختگی و تصادفی.} وقفه‌های اضافه، الگوی واقعی را با نوفه می‌پوشانند و دقت دسته‌بند را کم می‌کنند. اما نوفهٔ مستقل از قربانی را می‌توان با میانگین‌گیری از ردهای تکراری و فیلتر کردن تا حدی حذف کرد. برای اینکه الگو واقعاً پنهان شود، نوفه باید هم‌اندازهٔ سیگنال و از نظر شکل و مدت شبیه وقفه‌های واقعی باشد. هزینه: هدر دادن \lr{CPU} و انرژی و افزایش تأخیر. \textbf{میزان جلوگیری: متوسط، با هزینهٔ زیاد.}

\begin{center}\small
\small
\begin{tabular}{@{}>{\raggedright\arraybackslash}p{3.4cm}>{\raggedright\arraybackslash}p{4.4cm}>{\raggedright\arraybackslash}p{4.2cm}>{\raggedright\arraybackslash}p{2.3cm}@{}}
\toprule
  {\textbf{روش}} &
  {\textbf{مهم‌ترین نقطهٔ قوت}} &
  {\textbf{مهم‌ترین نقطهٔ ضعف}} &
  {\textbf{میزان جلوگیری}} \\
\midrule
  {۱) انتقال وقفه‌ها} &
  {حذف وقفه‌های دستگاه‌ها} &
  {زمان‌سنج، بین‌هسته‌ای، کش} &
  {متوسط} \\
  {۲) ساعت کم‌دقت} &
  {ساده و ارزان} &
  {حمله به ساعت دقیق نیاز ندارد} &
  {کم} \\
  {۳) وقفهٔ ساختگی} &
  {پنهان‌سازی الگوی واقعی} &
  {میانگین‌گیری و هزینهٔ زیاد} &
  {متوسط} \\
\bottomrule
\end{tabular}
\end{center}

\textbf{چرا وقفهٔ زمان‌سنج و وقفه‌های بین‌هسته‌ای را نمی‌توان به هستهٔ دیگر منتقل کرد؟}
\begin{itemize}
  \item \textbf{وقفهٔ زمان‌سنج:} هر هسته زمان‌سنج محلی خودش را دارد (مثلاً \lr{LAPIC timer}) و سیستم‌عامل برای زمان‌بندی روی همان هسته به آن نیاز دارد. سیستم‌عامل پیش از اجرای پردازه زمان‌سنج را تنظیم می‌کند تا پس از پایان سهم زمانی کنترل را پس بگیرد. اگر این وقفه به هستهٔ دیگر برود، هستهٔ اجراکنندهٔ پردازه هرگز پیش‌دستی نمی‌شود و سهم زمانی و زمان‌بندی از بین می‌رود.
  \item \textbf{وقفه‌های بین‌هسته‌ای (\lr{IPI}):} ماهیتاً پیامی از یک هسته به هستهٔ هدف مشخص است تا حالت همان هسته عوض شود؛ مثل ابطال \lr{TLB} همان هسته، بیدار کردن یا پیش‌دستی روی آن، یا اجرای یک تابع روی آن. منتقل کردن آن به هسته‌ای دیگر معنایش را از بین می‌برد. در مقابل، وقفهٔ دستگاه از کنترل‌گر وقفه به هر هسته‌ای قابل مسیریابی است.
\end{itemize}

